% !TEX root = Stolpe_2018_Thesis.tex
% !TEX encoding = UTF-8


\chapter{Related Work}
\label{cha:relatedwork}

%As tooling support for programming languages is a wide field of research, we only focus on work related to language implementation frameworks for dynamic languages and language-agnostic tooling approaches for \acp{ide}.

Our approach focuses on enabling performance-focused language implementation frameworks to provide their implemented dynamic languages with IDE-like tooling support with minimal effort. In this chapter, we describe related approaches and solutions. Furthermore, we compare them with our approach and implementation of \graallsp. As a major part of our approach is based on collecting and using run-time information to enhance certain \ac{lsp} features, we will also discuss work related to that field.

% To understand the challenging task of developing new languages and corresponding tools, we have to have a look at how existing and widely used \acp{ide} provide support for new languages. Among other \acs{ide}, Eclipse \cite{des2004eclipse}, VisualStudio \cite{Microsoft2018}, and IntelliJ \cite{JetBrains2018} provide extension points to develop \ac{ide}-specific plugins for new languages. Depending on the \ac{ide}, you have to implement lexer and parser yourself or provide at least grammar files to generate these components. Then, you have to manually implement further \ac{ide}-specific interfaces to support features like code completion, finding all references, or goto-definition. As already mentioned in \todo{motivation}, this will be difficult especially for dynamic languages. Another barrier for language or tool developers might be, that you have to implement these components in the \ac{ide}'s host language. Furthermore, if you want to support another \ac{ide}, the whole process needs to be repeated. The \ac{lsp} solves  some of these problems and the mentioned \acp{ide} support this protocol already. However, you still have to develope a language server for your language from scratch.

\ac{ide}-frameworks like Eclipse \cite{des2004eclipse}, Visual Studio\footnote{\url{https://visualstudio.microsoft.com} (accessed 2018-12-17)}, or IntelliJ\footnote{\url{https://www.jetbrains.com/idea} (accessed 2018-12-17)} provide extension points to develop \ac{ide}-specific plugins for new languages or can be used as language clients if they implement the \ac{lsp}. But as they do not support language developers in implementing programming languages in the first place and writing \ac{ide}-specific plugins or a language server still requires a lot of work, we do not further discuss them here. Furthermore, there are lots of language- and editor-specific tooling solutions for dynamic languages, which are also not our focus.

% \clearpage

% \paragraph{Outline}
% \begin{itemize}
%     % \item ??? Typical IDE-frameworks used by performance-focused language implementations
%     % \begin{itemize}
%     %     % Eclipse and IntelliJ -> you need to implement the custom IDE lang-specific parts in Java - you cannot use the language you implemented your language with
%     %     % You need to implement lexer and parser again and all other required interfaces - gap between impls
%     %     % These IDEs say, that you "only" have to implement the lang specific parts "with a small amount of effort"
%     %     % Features like code completion of properties and goto-definition are still hard to implement for dynamic languages
%     %     \item Eclipse
%     %     \item Visual Studio % https://docs.microsoft.com/en-us/azure/devops/integrate/ide/extensions/develop_extensions/add_new_language?view=vsts
%     %     \item IntelliJ 
%     %     % http://www.jetbrains.org/intellij/sdk/docs/reference_guide/custom_language_support.html
%     %     % http://www.jetbrains.org/intellij/sdk/docs/reference_guide/custom_language_support/code_completion.html
%     % \end{itemize}
%     \item Tools-first language implementation frameworks
%     \begin{itemize}
%         \item Language workbenches
%         \begin{itemize}
%             \item Xtext
%             % language development framework -> DSLs
%             % Xtext is a framework for development of programming languages and domain-specific languages. With Xtext you define your language using a powerful grammar language. As a result you get a full infrastructure, including parser, linker, typechecker, compiler as well as editing support for Eclipse, any editor that supports the Language Server Protocol and your favorite web browser
%             \item Spoofax and DynSem - generates a Truffle AST interpreter
%         \end{itemize}
%         \item Helvetia and Squimera
%     \end{itemize}
%     \item Language-agnostic Approaches Enabling \ac{ide}-like Features
%     \begin{itemize}
%         \item CTags
%         \item Kythe
%         \item srclib %https://srclib.org/
%         \item Lucid Energize 
%     \end{itemize}
%     \item Language-specific \ac{ide}-independent Feature Support
%     \begin{itemize}
%         \item Jedi (Python)
%         \item Rsense (ruby)
%         \item Tern (JS)
%     \end{itemize}
%     \item Improving \acs{ide} Features With Run-time Information
%     \begin{itemize}
%         \item Type harvesting
%         \item PyCharm - dynamic type inference for code completion
%         \item pyannotate - Insert annotations into your source code based on call arguments and return types observed at run-time
%         \item Hermion
%         \item Enhancing static ref search with dynamic data
%     \end{itemize}
% \end{itemize}


% \clearpage

\section{Tools-first Language Implementation Frameworks}
    \paragraph{Language Workbenches}
    Language workbenches are tools that support the efficient definition, reuse and composition of languages, especially \acp{dsl} \cite{Fowler2010DomainSpecificLanguages}, and their \acp{ide} \cite{erdweg2013state}. Hence, they provide the means to make the development and composition of new languages affordable, which has been proven by the usage of language workbenches and \acp{dsl} in various use cases in industry. Most language workbenches support a textual language notation, which can be edited in a free-form text editor. But there are also solutions like MPS \cite{voelter2012language}, which allow projectional editing or mathematical symbols. Furthermore, syntactic services (syntax highlighting, folding, or syntactic completion), semantic services (reference resolution, refactoring, diagnostics, or semantic completion), and further \ac{ide}-like features can be provided based on the language definitions.

    However, as already described in \autoref{cha:introduction}, these frameworks do not focus on the performance of the language implementations in the first place and are often bound to a limited set of \acp{ide}. Furthermore, dynamic languages cannot make use of all the services of language workbenches, as these services typically rely on the static information derived from the language definitions.

    Xtext \cite{eysholdt2010xtext} is one of the first language workbenches, which supports the \ac{lsp}, and which also proposes, that it can be used to develope \acp{dsl} and full-featured programming languages. However, Xtext faces the same challenges like the other workbenches, when it comes to dynamic languages, which we try to tackle explicitly in \graallsp.

    Visser et al. described an approach using the Spoofax language workbench \cite{kats2010spoofax} and the DynSem language \cite{vergu_et_al:LIPIcs:2015:5208} to systematically generate Truffle languages, i.e. \ac{ast} interpreters, from DynSem specifications \cite{visser2014language}. However, Spoofax provides plugins for the Eclipse \ac{ide} only. Supporting the \ac{lsp} seems to be planned, but is not yet present. Hence, it would be worth to evaluate, if a Truffle language generated through this approach could provide enough means to support at least some \graallsp~features. Nonetheless, the generated interpreters will still require manual tweaks, to benefit from all the GraalVM optimizations.
    %Language workbenches, a term popularized by Martin Fowler in 2005 [19], are tools that support the efficient definition, reuse and composition of languages and their IDEs. Language workbenches make the development of new languages affordable and, therefore, support a new quality of language engineering, where sets of syntactically and semantically integrated languages can be built with comparably little effort. This can lead to multi-paradigm and language-oriented programming environments [8, 61] that can address important software engineering challenges. \paragraph{Frameworks for DSLs}
    % \begin{itemize}
    %     \item not performance focused
    %     \item typically only one or a very limited set of target IDEs supported (e.g. an Eclipse plugin)
    %     \item not limited to textual notation (also graphical, tabular etc.)
    %     \item Composability
    %     \item Testing / debugging
    %     \item Editing mode: free form + projectional
    %     \item Syntactic services (Highlighting, Outline, Folding, Syntactic completion, Diff, Auto formatting)
    %     \item Semantic services (Reference resolution, Semantic completion, Refactoring, Error marking, Quick fixes, Origin tracking, Live Translation)
    % \end{itemize}
        
    %- stefan marr:  While Language Workbenches provide such features, I personally prefer the bottom up approach. Instead of specifying my language, I’d rather express the one truth of its semantics as an interpreter. In the end, I want to execute programs. So, let’s start there.

    % \paragraph{Xtext}
    % Xtext, part of the Eclipse project, is a framework for defining external DSLs in close cooperation with the Eclipse IDE. \cite{humer2015domain}

    % \paragraph{A Language Designer's Workbench} \cite{visser2014language} use Spoofax and DynSem \cite{vergu_et_al:LIPIcs:2015:5208}
    %     % To support the direct execution of programs in the object language, we systematically generate a Java-based AST interpreter from a DynSem specification for the language. Each constructor of the language is translated to a class in Java with an evaluate method based on the DynSem rules for the constructor.
    %     % DynSem specifications are executable through automatic generation of a Java-based AST interpreter.
    %     -> generated AST-interpreter for Truffle <- can one modify the generated interpreter?
    %     -> end-user IDE: bound to the Spoofax generated editor(-plugin)

    % \paragraph{Ensō} http://enso-lang.org
    %     -> Ensō is a self-describing workbench based entirely on interpreters implemented in an object-oriented language. Ensō avoids the generation of code and follows the hypothesizes, that interpreters are easier to develop and combine than compilers or transformational approaches. Because Ensō is implemented in itself, all available tools in the system can be reused to create applications. -> Smalltalk approach
    %     % Ensō rejects the common approaches for 1) model to code, 2) model to model, and 3) generic transformation languages.
    %     % Ensō is based entirely on interpreters not compilers/translators. Code generation is completely avoided, because it is complex and unnecessary. Ensō hypothesizes that interpreters are easier to develop, maintain, extend, and combine than compilers/transformations. As examples, both the Ensō data modeling language [1] and grammar language [2] are interpreted on the fly.
    %     % Ensō implements interpreters using ordinary object-oriented code. By defining interpreters as objects, traditional object-oriented composition and inheritance techniques can be used to extend and combine language semantics.
    %     % Ensō is based on external languages rather than embedding. An external language is one that has its own syntax and semantics, independent of any other language. An embedded language is defined as a library within a host language. Embedded languages typically reuse features of the host language, including variable binding or control flow.
    %     % Ensō strikes out on its own rather than building on existing standards, including UML and EMF.
    %     % Eclipse is a powerful and practical platform with many advantages. However, it is possible that much of its complexity is accidental rather than fundamental. If so, it may be possible to create an extensible integrated development environment with orders of magnitude less code and complexity than in Eclipse. To find out, Ensō will create its own IDE rather than build on top of Eclipse. Evidence from Newspeak, Cloud9, and Smalltalk suggest that this is possible.
    %     % Finally, Ensō is implemented in itself. This allows all tools written for the system to be reused in creating applications, and tools and libraries developed for applications to be reused in the creation of the system. It also means that the core languages and strategies that define the Ensō system can be modified to extend the functionality of Ensō. Smalltalk has shown the benefits of such deep reflection. Ensō provides a similar level of reflection, but in a more controlled way, as the system has a well-defined bootstrap process. Ensō is bootstrapped from Ruby, and is available today for experimentation.
    %     % [1] Alex Loh, Tijs van der Storm, William R. Cook. “Managed Data: Modular Strategies for Data Abstraction” To appear SPLASH Onward! 2012. http://www.cs.utexas.edu/~wcook/Drafts/2012/ensodata.pdf
    %     % [2] Tijs van der Storm, Alex Loh and William R. Cook. “Object Grammars: Compositional & Bidirectional Mapping Between Text and Graphs” To appear at the International Conference on Software Language Engineering (SLE), 2012. http://www.cs.utexas.edu/~wcook/Drafts/2012/ensogrammars.pdf
    %     % https://modeling-languages.com/enso-dont-design-your-programs-program-your-designs/
    
    \paragraph{Helvetia and Squimera}
    Two approaches, which focus on reusing existing tools of a live programming environment and the overall programming experience, are the Smalltalk-based \acp{ide} Helvetia \cite{renggli2010embedding} and Squimera \cite{niephaus2018live}. Helvetia enables embedding \acp{dsl} into Smalltalk, by adding hooks into the Smalltalk compiler chain. This way, the \acp{dsl} can be translated on demand into Smalltalk code, so that existing tools of the environment still work, without the need to generate and start a new \ac{ide}.
    
    Squimera is not a language implementation framework, but it enables polyglot programming and debugging through the composition of interpreters written in RPython. With some adjustments, some of the existing Smalltalk tools have been made reusable for writing and debugging Python and Ruby code, providing a consistent programming experience across languages. Like \graallsp, Squimera utilizes the interpreters and \acp{api} of a performance-focused language implementation framework, to retrieve data needed for the tools of the language users.
    %Like in Helvetia, no new \ac{ide} instance needs to be started to use other languages, as the environment is part of the Smalltalk language and runs in the same process.

    Helvetia and Squimera provide syntax highlighting and debugging support for the integrated languages, but features like code completion or goto-definition are only provided at Smalltalk level. However, the integrated languages could theoretically make use of all the tools of the Smalltalk environment. In Squimera, customized tools for polyglot programming of high-level languages have been created and are still a subject of research. This would not be possible in \acp{ide} using \graallsp~yet, as they are bound on the one hand to the features made available through the Truffle instrumentation \ac{api} and on the other hand to the specification of the \ac{lsp}.
    
    \paragraph{}
    In summary, language workbenches and the other described approaches can provide a richer, more customized, and domain-specific tooling support for new or existing languages than \graallsp. However, these solutions are not able to create high-performance dynamic languages together with tools for the users of the languages, supporting \ac{ide} features like code completion or goto-definition. Furthermore, the users are often bound to a single \ac{ide}, whereas \graallsp~allows them to choose from a variety of \ac{lsp} clients.


\section{Language-agnostic Approaches Enabling \ac{ide}-like Features}
    To the best of our knowledge, there are no approaches comparable to the \acl{lsp} yet, which are supported by almost all popular source code editors and languages. However, there exist other approaches for language-agnostic tooling support, which can be seen as predecessors or influencer for \ac{lsp}.

    \paragraph{Ctags}
    % \begin{itemize}
    %     % The ctags and etags programs (hereinafter collectively referred to as ctags, except where distinguished) generate an index (or "tag") file for a variety of language objects found in file(s). This tag file allows these items to be quickly and easily located by a text editor or other utility. A "tag" signifies a language object for which an index entry is available (or, alternatively, the index entry created for that object). Alternatively, ctags can generate a cross reference file which lists, in human readable form, information about the various source objects found in a set of language files. Tag index files are supported by numerous editors, which allow the user to locate the object associated with a name appearing in a source file and jump to the file and line which defines the name.
    %     % Exuberant CTAGS http://ctags.sourceforge.net/
    %     %\item \cite{hiebert1999exuberant} % Exuberant CTAGS http://ctags.sourceforge.net/
    %     \item \footnote{\url{http://ctags.sourceforge.net} (accessed 2018-12-17)}
    %     \item tagging a language object in an index file
    %     \item tag-file supported by editors
    %     \item allows user to locate the tagged object via name (goto-definition, list symbols)
    %     \item for new languages you need to define a regex or parser to find language objects
    %     \item only static analysis, no scoped completion (everything is suggested)
    % \end{itemize}
    Ctags is an old but still popular tool to create language-agnostic index (tag) files for symbols found in source code of various programming languages. The tags can be consumed by source code editors to provide a goto-definition search or a simple symbol completion. Over the years different flavors of the original Ctags have come up which extended the original file format for tags. A popular variant, Exuberant Ctags\footnote{\url{http://ctags.sourceforge.net} (accessed 2018-12-16)}, was originally distributed with Vim, but became a separate project, supporting over 40 programming languages. New languages can be added either by defining regular expressions to determine language symbols, or providing a parser.

    
    \paragraph{Energize}
    Energize/Cadillac \cite{gabriel1990foundation} was an object-oriented \ac{ide} for C++ built upon Emacs in the late 1980s. It was based on a language-agnostic client-server architecture for programming language environments. Clients (compilers, editors, debuggers, etc.) communicate with a server using protocols for a certain class of tools (compiler protocol, editor protocol, debugger protocol, etc.). The server is able to persistently store objects that represent source code, language elements, compiled code objects, and annotations (objects which are links between other objects), so that each tool can access the needed information. For example, annotations could be used to link compiler diagnostics to source code and/or language elements, but also to retain dependencies to enable incremental compilation. The persistent object store was intended to store information, such as dependency analyses, for the lifetime of the software. New languages could be added, by providing compilers for these languages, which support the compiler protocol.
    
    The \ac{lsp} could be interpreted as an evolved version of selected parts of the Energize protocols, especially of editor and compiler.
   
    % In the late 1980s, we started working at Lucid on a programming environment for C/C++. It was based on a programming-language-neutral, non-abstract-syntax-tree way of describing programs, annotations (a generalization of hypertext), a persistent repository for program information, and a tool-neutral user-interaction mechanism. To this day I have not seen a programming environment as advanced as Energize/Cadillac (official name/code name). There was really only one substantial paper published about the work, and that was in an obscure conference called C++ At Work (sorry if this offends someone). The paper was called Foundation for a C++ Environment (http://www.dreamsongs.com/Cadillac.html)

    \paragraph{Kythe}
    Kythe\footnote{\url{https://kythe.io} (accessed 2018-12-16)} is a pluggable, language-agnostic ecosystem for building tools that work with code. Kythe allows language analysis and indexing to be run as services, by defining language-agnostic protocols and data formats for representing, accessing, and querying source code information as data. It was originally developed as an internal Google project, due to the observation that most companies have projects which are created with more than three programming languages, but lack appropriate cross-language tools\footnote{\url{https://bsumm.net/2012/08/11/steve-yegge-and-grok.html} (accessed 2018-12-16)}. To support a new language in Kythe, a language compiler or interpreter needs to be extended, so that it can emit a rich \ac{ast}, symbol table and type graph in a Kythe-compliant intermediate representation. For dynamic languages however, the generated \acp{ast} do not include the information required by Kythe, so that additional type inference engines are needed.

    \paragraph{srclib}
    Another project which tried to tackle similar challenges as \ac{lsp} is srclib\footnote{\url{https://srclib.org} (accessed 2018-12-16)}, a polyglot code analysis library for static analysis and code search. srclib was the foundation of Sourcegraph\footnote{\url{https://sourcegraph.com} (accessed 2018-12-16)}, a code search and intelligence server that helps developers find, review, and understand code. Like \ac{lsp}, srclib defined a language-agnostic format, which can be consumed by development tools. New languages can be supported by implementing a tool chain that outputs the srclib format. For languages like JavaScript, Ruby, and Python, tool chains have been built reusing existing language-specific projects like Tern, YARD\footnote{\url{https://yardoc.org} (accessed 2018-12-16)}, and Jedi. However, with the upcoming of \ac{lsp}, the development of srclib was stopped and Sourcegraph also makes use of \ac{lsp} now.
    %-> not under active development anymore, Sourcegraph uses LSP now
    % srclib makes developer tools like code search and static analyzers better. It supports things like jump to definition, find usages, type inference, and documentation generation.
    % srclib consists of language analysis toolchains (currently for Go, Python, JavaScript, and Ruby) with a common output format, and developer tools that consume this format.
    % srclib originated inside Sourcegraph, where it powers intelligent code search over hundreds of thousands of projects.
    % srclib takes care of the code analysis that underlies key assistive features such as finding usages and jumping to definitions. It exposes a command-line API and outputs data in a well defined, language-independent format.
    % A plugin (or other tool) built on srclib just has to consume this data format to support all languages that srclib supports. srclib is designed to be modular and extensible. New languages can be added by implementing language toolchains that output the srclib format.
    % srclib is 100% open source and built on other well-established open-source projects like Tern, YARD, and Jedi.
    % Sourcegraph https://sourcegraph.com is a free, open-source, self-hosted code search and intelligence server that helps developers find, review, understand, and debug code. Use it with any Git code host for teams of any size. Start using it now, then install the Sourcegraph Docker image on your private code.
    % Solve problems before they exist, commit by commit. Code intelligence makes browsing code easier, with IDE-like hovers, go-to-definition, and find-references on your code, powered by Sourcegraph extensions and language servers based on the open-source Language Server Protocol

    \paragraph{}
    Like the \ac{lsp}, the described tools and approaches define more or less complex protocols to decouple languages and their development tools. Supporting a new language ranges from adding only a regular expression to implementing or instrumenting complex compilers, to output language-specific information in a language-agnostic format. However, code completion is not part of these protocols, which we wanted to support in \graallsp. Furthermore, besides \ac{lsp}, only Ctags is supported by a variety of source code editors and \acp{ide} at the date of writing.


\section{Language-specific Editor-agnostic Approaches Enabling \ac{ide}-like Features}
    As described in the introductory words to this chapter, there are lots of \ac{ide}- or editor-specific solutions to support code completion etc. for certain dynamic languages, which are not discussed here. But there are some language-specific solutions worth mentioning, which are designed to be editor-independent, like Jedi\footnote{\url{https://github.com/davidhalter/jedi} (accessed 2018-12-19)} for Python, Tern\footnote{\url{http://ternjs.net} (accessed 2018-12-19)} for JavaScript, or Rsense\footnote{\url{https://github.com/rsense/rsense} (accessed 2018-12-19)} for Ruby. Like \ac{lsp}, Tern and Rsense define a lightweight interchange format, so that editor-plugins can be created, which only have to map the requests and responses to their editor-specific \ac{api}. Tern and Rsense provide a server which accepts and responds JSON messages. However, the structure of the messages is influenced by the corresponding language and is therefore more specific than the message format of \ac{lsp}. Jedi provides a simple Python \ac{api} only. Therefore, editor-plugins need to be able to execute Python code or support the \ac{lsp}, because there already exist Python language servers using Jedi.

    All the described tools mainly rely on static code analysis, so that they all face limitations due to certain dynamic features of their corresponding language. Jedi was originally intended to provide a rich code completion by executing non-dangerous code, which did not work, so that a highly recursive static analyzer has been created\footnote{\url{https://jedi.readthedocs.io/en/latest/docs/features.html\#a-little-history}\\(accessed 2018-12-19)}. \graallsp~is able to execute source code to some extent yet, so that it supports code completion also for dynamic features. However, as the described solutions are language-specific, they can support certain features more suitable and with more detail than a generic \ac{lsp}. But this comes at another cost, because these tools cannot reuse the existing language parsers to get the required data. They had to built and maintain a separate parser, so that when the languages evolve, there might always be a gap between language elements already supported in the original parser, which have not been implemented in the parsers of the tools yet. In contrast, \graallsp~can indirectly utilize the existing language parsers through the Truffle \ac{api}.


\section{Improving \acs{ide} Features With Run-time Information}
\label{sec:rw:collectionRunTimeInformation}
    % \begin{itemize}
    %     \item Type harvesting \cite{haupt2011type}
    %     \item Hermion \cite{rothlisberger2008exploiting}
    %     %\item Nemo Bachelor Project? -> sandboxed execution in Squeak
    %     \item https://github.com/dropbox/pyannotate ->  Insert annotations into your source code based on call arguments and return types observed at runtime. - uses Python's sys.setprofile(profilefunc)
    %     \item https://blog.jetbrains.com/pycharm/2013/02/dynamic-runtime-type-inference-in-pycharm-2-7/
    %     \item Enhancing Static Source Code Search With Dynamic Data \cite{holmes2010enhancing}
    % \end{itemize}

    A major part of our approach focuses on collecting and utilizing run-time information to answer \ac{lsp} requests. Therefore, we have a look at where and how other approaches apply run-time information to improve \ac{ide} features.

    \paragraph{Hermion}
    % We address these questions and present our working prototype, an IDE called Hermion with the capability to capture runtime information from an application under development and to exploit this information by enhancing navigation and browsing of the source code.
    The Hermion \ac{ide} \cite{rothlisberger2008exploiting} collects run-time information to enhance program comprehension by improving static source code navigation and browsing in Squeak/Smalltalk. In detail, they apply a filter to the list of methods of the implementors view to display only those methods which have been called in a specific scenario. Furthermore, they display collected type information for instance variables and provide an additional reference view in the method source code view, that shows a list of observed classes which are referenced in the current method.
    
    Hermion builds upon partial behavioral reflection \cite{tanter2003partial} to restrict the collection of run-time information to certain source code parts and detail levels. However, due to performance issues, the users of the Hermion \ac{ide} have to make a trade-off between performance and the amount of collected information. The users need to manually specify which classes shall be instrumented, to keep the cost of instrumentation and reflective operations at an acceptable level. In \graallsp, the set of \ac{ast} nodes which shall be instrumented when collecting and caching run-time information, is implicitly defined by the \ac{lsp} workspace, i.e. all parsed source code files. As we collect snapshots of frames, we increase performance and reduce the amount of collected instances by taking only the first frame which is reported for every instrumented node. As Hermion continuously observes the instrumented classes and can for example collect different types for a variable, which \graallsp~would not be able to detect.
    
    % \begin{itemize}
    %     \item Hermion builds upon partial behavioral reflection \cite{tanter2003partial} to restrict the collection of run-time information only to certain source code parts and detail levels. This is similar to Truffle's instrumentation of relevant \ac{ast} nodes, which we use in \graallsp. % E. Tanter, J. Noye, D. Caromel, and P. Cointe. Partial Behavioral Reflection: Spatial and Temporal Selection of Reification. In Proc. of the Conference on Object-oriented Programing, Systems, Languages, and Applications, pages 27–46. ACM, October 2003
    %     % Another approach to gain performance is partial behavioral reflection [Tanter et al. 2003], which restricts the application of a MOP in spatial and temporal manner. Instead of applying to all program parts equally, the developer optimizes the MOP’s application to only the places where it is needed to avoid unnecessary overhead. However, this does not reduce the cost of reflective operations used. Furthermore, it burdens again the programmer with the optimization while dispatch chains are part of the language implementation and thus optimize reflective operations without programmer intervention
    %     \item developer needs to manually specify which classes shall be instrumented
    %     \item due to performance issues, developers have to make a trade-off between performance and collected information
    %     \item similar to us: they rely on coverage of executed code
    % \end{itemize}

    \paragraph{Enhancing Static Source Code Search in Eclipse}
    Holmes and Notkin \cite{holmes2010enhancing} described an approach to improve the static find-references feature of the Eclipse \ac{ide} using collected run-time information. Only executed methods are included in the search results, so that users benefit from a limited result set related to their current application. To collect the required run-time information, they use a custom AspectJ \cite{kiczales2001overview} aspect for the application's test harness. % Holmes and Notkin [15] describe an approach when the “find references” capability of an IDE is filtered using information from dynamic analysis – only methods executed in the given scenario are returned.

    Our \graallsp~prototype does not utilize run-time information to answer the \emph{references} request of the \ac{lsp} yet. We perform a static symbol-based search in files of the current workspace, but it should be possible to enhance our implementation of the \emph{references} request in a similar way.

    \paragraph{PyCharm Debugger}
    PyCharm is a feature-rich \ac{ide} for the Python programming language. The debugger of PyCharm is able to optionally collect run-time type information during a debug session\footnote{\url{https://blog.jetbrains.com/pycharm/2013/02/dynamic-runtime-type-inference-in-pycharm-2-7} (accessed 2018-12-16)}. These information are used to improve code completion and documentation popups during and after the debugging session. The collected type information can also be persisted as docstrings in Python source code.
    % PyCharm’s debugger is now capable of collecting type information at runtime and using this information for improving code completion suggestions, code insight and quick documentation popups. This needs to be enabled in the debugger settings.
    % You debug your code all the time, and now in the course of debugging you can also collect type information. Afterwards, it will be used in auto-completion, and you will be able to put the collected type information into docstrings as well. Please note that enabling this feature may affect debugger performance, so use it with care.

    Similar to our approach, the PyCharm's debugger can utilize unit tests as entry point. However, \graallsp~does not require to run code in a special debug mode to collect run-time information and can also make use of these information to provide a precise goto-definition for callables. For the PyCharm \ac{ide}, this kind of dynamic run-time type inference is a small addition to the other mature features, which rely on static analysis, whereas for \graallsp, most features rely on run-time information. Like the previously described Hermion \ac{ide}, PyCharm's debugger collects types and is able to remember all type variants of variables, whereas \graallsp~collects a frame snapshot per instrumented node. This way, \graallsp~can provide code completion items also for properties of an object, which are not type related.


    \paragraph{}
    % \begin{itemize}
    %     \item all these solutions are language-specifically implemented
    %     \item focus on a single aspects (code navigation, information on hover, or code completion)
    %     \item they need to be triggered manually, we do auto-exec in sandbox
    %     \item we can do a specific run to the location we are interested and collect only a single run-time object
    % \end{itemize}
    In summary, all the three described solutions collect language-specific run-time information to improve certain static aspects of existing language-specific tooling. Hermion focuses on code navigation and comprehension, the approach of Holmes and Notkin on finding all references, and PyCharm on code completion and documentation. In contrast, \graallsp~collects run-time information in a language-agnostic way and tries to support as much \ac{ide} features as possible with these information. \graallsp~uses static approaches rather as fallback when no run-time information can be retrieved.

    \paragraph{}
    The following two approaches do not directly target to improve specific \ac{ide} features, but they describe how to collect run-time information, so that other tools can make use of them. 

    \paragraph{PyAnnotate}
    PyAnnotate\footnote{\url{https://github.com/dropbox/pyannotate} (accessed 2018-12-16)} is a Python library to collect type information (call arguments and return types) at run-time, using a Python-specific profiling hook. After a collection run, these information can be automatically inserted as annotations or comments into the Python source code, so that other tools can make use of them.

    \graallsp~collects run-time objects to answer \ac{lsp} requests only. The objects or their types are not persisted in the source code files for the later use of other tools.

    \paragraph{Type Harvesting}
    Haupt et al. have presented an approach called \emph{type harvesting} to collect type information for dynamic programming languages \cite{haupt2011type}. They suggest to use these information for code generation, optimization, or program comprehension. Their solution is based on a stepwise execution of bytecode while observing the types of relevant components. The set of relevant components, which shall be observed, are called an \emph{acre}, which have to be determined manually by the user before the harvesting is started. A \emph{bucket} stores mappings of bytecode instructions, arguments, variables, etc., to the collected concrete types. Like in our approach, they suggest to use unit tests as entry points to start code execution.

    In contrast to our approach, \emph{type harvesting} continuously observes the stepwise execution of bytecodes, leaving the actual bytecode instruction stream untouched. However, the simulated stepwise execution is quite slow and they need to check for every step, if the current bytecode is on an acre. \graallsp~benefits from Truffle's instrumentation and optimizations of \ac{ast} nodes, so that the execution of instrumented nodes has a very small overhead.

    Another difference is, that type harvesting stores specific mappings of bytecode instructions to types, whereas \graallsp~stores lexical mappings of lines and columns to frames. Storing types only is obviously more memory efficient, but would not work for \graallsp, as we need the actual run-time objects to utilize the Truffle \ac{api}. 


%%% Local Variables:
%%% mode: latex
%%% TeX-master: "Stolpe_2018_Thesis"
%%% coding: utf-8
%%% End: